\subsection{Compiler Basics}

A \emph{compiler} is commonly understood as a program that converts a high-level language to a low-level language executable by hardware.
However, in general a compiler translates one representation of a program to another, and perhaps does some analysis along the way. 

All translation is transformation. A Greek-to-English 
translation of the Odyssey will preserve the meaning
while reading as close as possible to the original. Code translation is 
the same. There are many possible translations of a given program 
to another language, with different properties and optimizations. 
Two translations that accomplish the same task may run in different times or
use different amounts of memory. 

Compilers save poor programmers from writing directly in assembly.
The many flavours of high-to-low compilers enable us to write in one language and compile 
to many different architectures, instead of needing to write 
the same program in x86 for x86 machines, RISC-V for RISC-V chips, and so on. 

